

\section*{Conclusion}
\label{sec:conclusion}

In this manuscript, we introduce \dolma, a three trillion token English corpus for language model pretraining. 
The \dolma corpus is comprised of a diverse set of content, including web documents, scientific papers, code, public-domain books, social media, and encyclopedic materials.
Building off a list of explicit desiderata, we document our data curation pipelines, providing experimental results that support our decisions. 
We freely release \dolma and open-source all tools we used to curate this dataset as part of the OLMo project~\citep{olmo20247b}. 
Since the time of writing, we have made improvements to \dolma and have continued to make releases; for example, our follow-on release of \textbf{\dolma~\texttt{v.1.7}} yields significant performance improvement on downstream tasks, holding the model constant.\footnote{
\href{https://blog.allenai.org/olmo-1-7-7b-a-24-point-improvement-on-mmlu-92b43f7d269d}{\path{medium.com/p/92b43f7d269d}}}
We hope this line of work can promote transparency, reproducibility, and further research in the field of language modeling, 
as well as address the current gap in the availability of pretraining data of commercial and open language models.
We release \dolma under ODC-By
and our toolkit under Apache 2.0.
